%% PLANTILLA PARA ARTICULOS

\documentclass[10pt,twoside,spanish,a4paper]{article}
%%%%%%%%%%%%%%%%%%%%%%%%%%%%% no tocar
% Alfredo Cuesta. 01/01/2013. acuestai@ucm.es

%si se quiere usar ARIAL hay que poner estas 2 lineas:
%\usepackage[scaled]{uarial} %DESCOMENTAR
\renewcommand*\familydefault{\sfdefault}
%

\usepackage[compact]{titlesec}
 \titlespacing\section{0pt}{12pt plus 0pt minus 0pt}{0pt plus 0pt minus 0pt}
 \titlespacing\subsection{0pt}{5pt plus 0pt minus 0pt}{0pt plus 0pt minus 0pt}
 \titlespacing\subsubsection{0pt}{0pt plus 0pt minus 0pt}{0pt plus 0pt minus 0pt}


%\usepackage[T1]{fontenc}
\usepackage[utf8]{inputenc}
%\usepackage[latin9]{inputenc}
\usepackage{
amssymb,
graphicx,
color,
url
%wrapfig %Inclusión de gr·ficos al lado de texto
}

\linespread{1.3}%separacion de lineas
\parindent=0.0in%indentancion de parrafo
\parskip=6pt%separacion entre el parrafo siguiente
\setlength{\topmargin}{-0.5in}
\setlength{\oddsidemargin}{0in}
\setlength{\evensidemargin}{0in}
\setlength{\textwidth}{16cm}%{4.72in}
\setlength{\textheight}{22cm}%{7.08in}

\usepackage{babel}
\addto\shorthandsspanish{\spanishdeactivate{~<>}}
% correct bad hyphenation here
\hyphenation{nu-es-tro co-no-ci-mien-to par-ti-cu-lar ca-rac-te-ris-ti-cas ni-vel se-gui-mien-to que-re-mos ca-ta-lo-go}
\date{}


% --- customization ---

\renewcommand\labelitemi{$\bullet$}
\let\olditemize\itemize
\renewcommand{\itemize}{
  \olditemize
  \vspace{-0pt}
  \setlength{\itemsep}{6pt}
  \setlength{\parskip}{0pt}
  \setlength{\parsep}{0pt}
  %\setlength{\topsep}{-15pt}
}

\let\oldenumerate\enumerate
\renewcommand{\enumerate}{
  \oldenumerate
  \vspace{-2pt}
  \setlength{\itemsep}{1pt}
  \setlength{\parskip}{0pt}
  \setlength{\parsep}{0pt}
  %\setlength{\topsep}{-15pt}
}

\usepackage{verbatim}
\usepackage{fancyhdr}
\usepackage{multirow}
\usepackage{colortbl}
\usepackage{verbatim}
\usepackage{amsmath}

\definecolor{olivegreen}{rgb}{0.0,0.5,0.0}

\begin{document}

\renewcommand{\tablename}{Tabla}
\renewcommand{\figurename}{Figura}
\title{
  \vspace{-0.2cm}
	\begin{center}
		\large
		\begin{tabular}{rrl}
		\multirow{4}{*}{\includegraphics[width=1.7cm]{UCM.eps} \hspace{1cm}}&Grado: & Ingeniería Electrónica de Comunicaciones\\
		&Asignatura:& Procesamiento de Señales \\
		&Profesores:  & Jesús Chacón Sombría \\
		&Profesores:  & María José Gómez Silva \\
		\end{tabular}
	\end{center}
	\rule{\textwidth}{0.4pt}
	%
	\vspace{-1.5cm}
	\begin{center}
		\color{blue} \Huge
		\begin{tabular}{c}
	        \textbf{Ejercicios}\\
		\textbf{Muestreo y Reconstrucción}\\
		\end{tabular}
	\end{center}
	\vspace{-0.5cm}
	\rule{\textwidth}{0.4pt}
}
\maketitle
 \thispagestyle{empty}
\pagestyle{fancy}

\lfoot{\includegraphics[width=0.8cm]{UCM.eps}}
\cfoot{\vspace{-0.6cm} Ingeniería Electrónica de Comunicaciones }
\rfoot{\vspace{-0.6cm} \thepage}
\lhead{Práctica}
\rhead{Práctica 2A: Variables aleatorias}
\renewcommand{\headrule}{ \vspace{-0.2cm} \rule{\textwidth}{0.4pt}}
\renewcommand{\footrule}{\vspace{0.5cm} \rule{\textwidth}{0.4pt}\\ \vspace{-0.3cm} \\}

%**************************************************************<- para indicar pregunta de la pr·ctica
\newcounter{ctPreg}
\setcounter{ctPreg}{1}
\newcommand{\itemP}{ \item[{\color{blue}\fbox{\textbf{Ejercicio \arabic{ctPreg}:}}}] \addtocounter{ctPreg}{1}}  %
\newcommand{\itemB}{ \item[{\color{blue}$\bullet$} ]    }
\newcommand{\itemR}{ \item[{\color{blue}$\triangleright$} ]    }       %
%**************************************************************

\vspace{-8em}

	\begin{itemize}
		\itemP La señal $x(t)$ tiene ancho de banda $\omega_m$. Determinar la frecuencia de muestreo mínima para que no se produzca aliasing:
		\begin{itemize}
			\item $x(t) + x(t-1)$
			\item $\frac{dx(t)}{dt}$
			\item $x^2(t)$
			\item $x(t)cos(\omega_0 t)$
		\end{itemize}

		\itemP Determinar el periodo de muestreo necesario para poder reconstruir las siguientes señales:
		\begin{itemize}
			\item $x_1(t)=\mathrm{e}^{-t}, t \geq 0$
			\item $x_2(t)=2f_c\mathrm{sinc}(2f_c t)$
			\item $x_3(t)=\mathrm{sinc}^2(2f_c t)$
		\end{itemize}

		\itemP Indica la respuesta impulso del retenedor de orden 1 (FOH) causal y del no causal. Compara sus espectros frecuenciales mediante
		la transformada de Fourier.

		\itemP Muestrea las siguientes señales, eligiendo un valor adecuado de $T_s$. Indica en cada caso si es posible o no 
		evitar el aliasing, y reconstruye la señal utilizando retenedores de orden cero y uno.
		\begin{itemize}
			\item $x_1(t)=\mathrm{sin}(\frac{2\pi}{4} t)$
			\item $x_2(t)=\mathrm{sin}(\frac{2\pi}{4} t), 0\leq t\leq 7$
			\item $x_3(t)=\mathrm{sin}(\frac{20\pi}{700}t)\mathrm{cos}(\frac{20\pi}{35}t)$
		\end{itemize}

		\itemP Un conversor A/D de 10 bits admite una entrada entre 0 y $V_{ref}$. Según la configuración, $V_{ref}$ puede tomar diferentes valores (1.1 V, 3.3 V, 5 V).
		\begin{itemize}
			\item Indica la resolución, número de niveles y rango dinámico para cada posible configuración.
			\item ¿Cuantos bits serían necesarios, en cada caso, para obtener un precisión de $1e^{-4}$ V?.
		\end{itemize}

	\end{itemize}

\end{document}